\section{反代数}\label{opposite}

在第 \ref{lagbr} 节我们解释了：若 $W$ 是复流形、典则丛为 $K$，且 $Y=T^*W$ 带标准复辛形式 $\Omega=\sum_i\d p_i\d q^i$，则 $\A=\Hom(\B_\cc,\B_\cc)$ 是作用在 $K^{1/2}$ 截面上的微分算符代数。这一陈述在限制到 $W$ 的开集后仍成立。

这里我们沿更标准的数学路线解释这一点。下面的构造与第 \ref{hermitian} 节中 Hermite 内积的构造有些相似，但我们使用的是线性而非反线性运算。

若 $\A$ 是代数，反代数 $\A^\op$ 定义如下：$\A^\op$ 的元素 $a^\op$ 与 $\A$ 的元素 $a$ 一一对应，但按相反次序相乘（故 $a^\op b^\op=(ba)^\op$）。

在 $Y=T^*W$ 的形变量子化情形，代数 $\A=\Hom(\B_\cc,\B_\cc)$ 同构于其自身的反代数。这可以如下看出。$Y$ 有全纯对合 $\tau:(p,q)\mapsto(-p,q)$，因而满足 $\tau^*(\Omega)=-\Omega$。在 $A$-模型中，可通过反转模型所定义二维流形 $\Sigma$ 的定向来补偿 $\Omega$ 符号的反转。因此 $\Sigma$ 上的 $A$-模型在 $\Sigma$ 的定向反转微分同胚连同靶空间上 $\tau$ 的作用下不变。这样的微分同胚反转 $\partial\Sigma$ 的定向，从而反转算符在 $\partial\Sigma$ 上插入的次序。但反转边界上算符插入的次序，等价于把 $\A$ 与 $\A^\op$ 互换。结论是：代数 $\A$ 通过一个在经典层面作用为 $(p,q)\to(-p^\op,q^\op)$ 的自同构，同构于其反代数 $\A^\op$。

另一方面，设 $\DD$ 是作用在全纯线丛 $\L$（或线丛的复幂）截面上的微分算符代数。我们马上会解释，反代数 $\DD^\op$ 是作用在 $K\otimes\L^{-1}$ 截面上的微分算符代数。因此 $\DD$ 同构于 $\DD^\op$ 意味着 $\L$ 同构于 $K\otimes\L^{-1}$，即 $\L=K^{1/2}$。于是，若 $\A$ 同构于 $\A^\op$，它必是作用在 $K^{1/2}$ 截面上的微分算符代数。

关于 $\DD^\op$ 的断言可用形式的分部积分证明。回忆在复流形上，外微分有展开 $\d=\partial+\bar\partial$，其中 $\partial$ 把微分形式的次数移动 $(1,0)$。典则丛 $K$ 就是 $(n,0)$-形式丛。对 $(n-1,0)$-形式 $\lambda$，$\partial\lambda$ 是 $(n,0)$-形式，即 $K$ 的截面。若 $D$ 是作用在 $\L$ 截面上的微分算符，其“转置” $D^\tr$ 是作用在 $K\otimes\L^{-1}$ 截面上的微分算符，定义如下。设 $g$ 是 $\L$ 的局部截面，$f$ 是 $K\otimes\L^{-1}$ 的局部截面。则 $D^\tr$ 由下述条件定义：对任意这样的 $f,g$，
\be\label{kilox}fDg=(D^\tr f)g+\partial\lambda,\ee
其中 $\lambda$ 是由 $f,g$ 及其导数局部构造的 $(n-1,0)$-形式。由此立即推出：若 $D_1,D_2$ 是作用在 $\L$ 截面上的两个微分算符，则 $(D_1D_2)^\tr=D_2^\tr D_1^\tr$。由于乘法次序反转，这表明对应 $D\to D^\tr$ 把 $\DD$ 映到其反代数 $\DD^\op$。因此，若 $\DD$ 是作用在 $\L$ 截面上的微分算符代数，则 $\DD^\op$ 是作用在 $K\otimes\L^{-1}$ 截面上的微分算符代数。

作为例子，设 $W$ 是 Riemann 面，$\L$ 是平凡线丛 $\O$，$D$ 是局部可写成 $D=a(z)\partial_z$ 的微分算符。设 $g$ 是 $\O$ 的局部截面，$f$ 是 $K$ 的局部截面。恒等式
\be\label{zor}fa\partial_zg=\left(-\partial_z(af)\right)g+\partial_z(fag)\ee
表明 $D^\tr$ 可局部定义为 $D^\tr f=-\partial_z(a(z)f(z))$。这里的负号反映了如下事实：在经典层面，从 $\DD$ 到 $\DD^\op$ 的映射来自作用为 $(p,q)\to(-p,q)$ 的对合。
